\documentclass[11pt,a4paper]{article}
\usepackage[utf8]{inputenc}
\usepackage[T1]{fontenc}
\usepackage{amsfonts}
\usepackage[french]{babel}
\frenchbsetup{StandardLists=true} % à inclure si on utilise \usepackage[francais]{babel}
\usepackage{enumitem}
\usepackage{amssymb}
\usepackage{amsmath}
\usepackage[left=2cm,right=2cm,top=2cm,bottom=2cm]{geometry}
\usepackage{multirow}
\usepackage[dvipsnames]{xcolor}

\usepackage[T1]{fontenc}
\usepackage{fouriernc}
\usepackage{graphicx}

\usepackage{setspace}
\singlespacing

\usepackage{pstricks}
\usepackage{xkeyval}
\usepackage{pst-labo}


\usepackage{multicol}
\usepackage{caption}
\usepackage{fancyhdr}
\pagestyle{fancy}
\renewcommand\headrulewidth{1pt}
\fancyhead[L]{Lycée Denis Diderot }
\fancyhead[R]{Classe de TSN1}
\setcounter{page}{10}\fancyfoot[C]{page \thepage}
\fancyfoot[R]{/today}
\fancyfoot[L]{Alexis RUIZ}
\usepackage{multirow,tabularx,booktabs}
\newcommand{\cadreblanc}[1]{%
\medskip\framebox[\linewidth][c]{%
\rule{0mm}{#1mm}}\par
\medskip}
\setlength{\columnseprule}{.25mm}

\usepackage{awesomebox}
\setlength{\aweboxvskip}{0mm}
\usepackage{pifont}
\usepackage{chemmacros}
\chemsetup{modules=all}

\renewcommand{\arraystretch}{2}
\newcommand*\acocher{\quitvmode{\fboxsep0pt \fboxrule0.8pt \fbox{\vrule height1.8ex depth.1ex width0pt \vrule height0pt depth0pt width1.9ex }}}



\newcommand{\defproblem}{\awesomebox[black]{2pt}{\rotatebox{90}{\large{\textbf{Problématique}}}}{black}}
\newcommand{\defconclusion}{\awesomebox[black]{2pt}{\rotatebox{90}{\Large{\textbf{Conclusion}}}}{black}}
\newcommand{\appel}{\awesomebox[black]{2pt}{
\includegraphics[scale=0.5]{images/appel exam.png}}{black}}
\newcolumntype{R}[1]{>{\raggedleft\arraybackslash }b{#1}}
\newcolumntype{L}[1]{>{\raggedright\arraybackslash }b{#1}}
\newcolumntype{C}[1]{>{\centering\arraybackslash }b{#1}}

\frenchbsetup{StandardLists=true}
\begin{document}



\center \LARGE{TP1 - Mesurer la longueur d'onde d'une onde sonore}
\\[0,3cm]

\normalsize

\hrule\
\\[0,2cm]

\flushleft \textbf{NOM :\\
Prénom:\\[0,2cm]
}
\hrule\
\\[0,2cm]


\ding{202} \text{Expérience :} \\
\begin{center}
    \includegraphics[scale=0.5]{Chapitre 4 - Propagation du son/images/montageUS.png} 
\end{center}

\begin{enumerate}
    \item Réaliser le montage (ne pas toucher aux réglages), on placera le récepteur à 10 cm devant et dans l'axe de l'émetteur.
    \item Mettre le générateur en marche puis régler l'oscilloscope pour obtenir au moins deux périodes des signaux observés. 
    \appel{\fbox{
\begin{minipage}{13.3cm}{
\begin{center}
 Appeler le professeur pour pour qu'il vérifie votre montage  et vos réglages.
 \end{center} }

\end{minipage}
}}

    \item Représenter les deux signaux :\\[5 pt]
\begin{center}
    
\def\scl{0.8}%scaling factor of the picture
\begin{tikzpicture}[
  scale=\scl,
  controlpanels/.style={yellow!30!,rounded corners,draw=black,thick},
  screen/.style={white,draw=black,thick},
  trace/.style={white!60!yellow!20!, ultra thick},
  smallbutton/.style={white,draw=black, thick},
  axes/.style={thick}]

  % Screen, centered around the origin then shifted for easy plotting
  \begin{scope}[xshift=7cm,yshift=8cm,samples=150]
    \fill[black!30!,rounded corners,draw=black,thick](-5.3,-4.3)
      rectangle (5.3,4.3);
    \fill[screen] (-5.0,-4.0) rectangle (5.0,4.0);
    \draw[thin] (-5.0,-4.0) grid (5.0,4.0);
    \draw[axes] (-5,0)--(5,0); % Time axis
    \draw[axes] (0,-4)--(0,4);
    \foreach \i in {-4.8,-4.6,...,4.8} \draw (\i,-0.1)--(\i,0.1);
    \foreach \i in {-3.8,-3.6,...,3.8} \draw (-0.1,\i)--(0.1,\i);
  \end{scope}

  
\end{tikzpicture}
\end{center} 
\item Que peut-on dire de la période du signal émis et de la période du signal reçu ?\\[2pt]
\cadreblanc{8}

\item Calculer, en seconde (s), la période des signaux \\[2pt]
\cadreblanc{8}

\item Calculer, en Hertz (Hz), la fréquence.\\[2pt]
\cadreblanc{8}

\item Pourquoi le son n'est pas audible ?\\[2pt]
\cadreblanc{8}



\item Déplacer le récepteur pour que les deux signaux soient en phase, on superposera les deux courbes, 

\appel{\fbox{
\begin{minipage}{13.3cm}{
\begin{center}
 Appeler le professeur pour pour qu'il vérifie votre réglage et vous explique comment réaliser la mesure suivante. 
 \end{center} }

\end{minipage}
}}

Noter la distance $d_1$ entre l'émetteur et le récepteur.~~~
\begin{center}
    \fbox{$$d_1=~~~~~~~~~~~~~~ cm$$}
\end{center}

\vspace{3mm}
\item Cette mesure est peu précise, nous allons donc répéter l'opération 10 fois puis obtenir une nouvelle distance $d$ qui sera plus précise. Noter, en cm, la valeur moyenne de $d$
\begin{center}
    \fbox{$$d =~~~~~~~~~~~~~~ cm$$}
\end{center}

\end{enumerate}

\vfill 

\ding{203} \textbf{Résultats} : \\
\begin{center}
    
\fbox{
\begin{minipage}{12cm}{
\begin{center}
    
\textbf{Cette distance s'appelle la longueur d'onde de l'onde sonore.\\
Elle se note $\lambda$ (lambda). \\
Elle se mesure en mètre.\\[10pt]
\LARGE{$\lambda = ~~~~~~~~~~ m$}
}
\end{center}}
\end{minipage}
}
\end{center}

\vfill 

\ding{204} \textbf{Calcul de la vitesse de l'onde sonore} : \\[2pt]
A l'aide des résultats précédents déterminer, en m/s, la valeur de la vitesse du son dans l'air :\\[2pt]
\cadreblanc{29}

\vfill 

\ding{205} \textbf{Applications} : \\[2pt]
Lors d'un orage le bruit du tonnerre accompagne les éclairs produits par la foudre.\\
Pourquoi entendons-nous le tonnerre après avoir vu l'éclair ? \\[10pt]
\dotfill

\vspace{10pt}

\dotfill

\vfill 

Comment déterminer la distance approximative à laquelle est tombée la foudre ? \\[10pt] 
\dotfill

\vspace{10pt}

\dotfill

\vfill
\end{document}